\documentclass[runningheads]{llncs}

\usepackage[T1]{fontenc}
\usepackage{amsmath}
\usepackage{amssymb}
\usepackage{listings}
\usepackage{url}

\lstset{
  basicstyle=\ttfamily\small,
  breaklines=true,
  columns=fullflexible,
  keepspaces=true,
  showstringspaces=false
}

\title{Recasting Classical Deductive Program Synthesis Examples in SNARK}
\titlerunning{Deductive Synthesis Examples in SNARK}

\author{Codex}
\authorrunning{Codex}
\institute{SNARK workspace notes, \url{https://github.com/nilqed/SNARK}}

\begin{document}

\maketitle

\begin{abstract}
SNARK supports answer extraction from refutation proofs, a mechanism that
connects automated theorem proving with deductive program synthesis.  This note
documents a small extension of the local SNARK example suite under
\texttt{examples/synthesis}.  The extension recasts two classical
Manna--Waldinger developments as runnable SNARK examples: synthesis of the
unification algorithm and emergence of the binary-search paradigm in the
square-root approximation problem.  The development is not a verbatim
reproduction of the original deductive-tableau proofs.  Instead, it isolates
the central constructive obligations and uses SNARK's conditional answers to
extract program terms from proofs.

\keywords{SNARK \and Answer extraction \and Deductive synthesis \and
Unification \and Binary search}
\end{abstract}

\section{Overview}

Deductive program synthesis treats program construction as theorem proving.
The desired program is specified by an existential theorem: for each admissible
input there exists an output satisfying the specification.  A constructive
proof of that theorem carries enough information to compute the output.  In
SNARK, this information is represented by answers attached to rows.  When a
proof closes, \texttt{(answer t)} can be read as the extracted witness term.

The local development adds a new synthesis example file.  It contributes five
entry points:
\begin{lstlisting}
(mw-unification-variable-case)
(mw-unification-constant-case)
(mw-unification-compound-case)
(mw-square-root-binary-search-step)
(mw-square-root-binary-search-program)
\end{lstlisting}
The aggregator \texttt{(mw-synthesis-examples)} runs all five and returns their
extracted program terms.  The examples are also wired into
\texttt{snark-system.lisp} and \texttt{snark-examples.asd}, so rebuilding SNARK
compiles them with the rest of the example suite.

\section{Answer Extraction in the Development}

Each example follows the same pattern.  First, it initializes SNARK with
resolution and conditional answer creation:
\begin{lstlisting}
(initialize)
(use-resolution t)
(use-conditional-answer-creation t)
\end{lstlisting}
Next, it declares an abstract object language of constructors, predicates, and
program symbols.  The specification cases are asserted as implications.  The
goal is then proved with an answer template:
\begin{lstlisting}
(prove '(some-output-relation input ?result)
       :answer '(values ?result))
\end{lstlisting}
If two proof branches contribute different answers, SNARK combines them with
\texttt{ANSWER-IF}.  This is the key mechanism used here to expose conditional
branches in the synthesized program.

\section{Working Example I: Unification}

Manna and Waldinger's unification paper specifies unification as the task of
finding a substitution that makes two expressions identical, or reporting
failure when no such substitution exists \cite{MannaWaldinger1981}.  Later
presentations emphasize the same existential-output view and its fit with
SNARK-style proof search \cite{Waldinger2021}.  The local SNARK example uses a
small first-order term language:
\begin{lstlisting}
mw-var, mw-const, mw-app        ; expression constructors
mw-bind, mw-compose            ; substitution constructors
mw-unify, mw-subst, mw-subst-of
mw-none, mw-some               ; failure/success result constructors
\end{lstlisting}

\subsection{Variable Case}

The variable case encodes the occurs check:
\begin{align*}
  \mathrm{occurs}(\mathrm{var}(v),t)
    &\Rightarrow \mathrm{unify\_out}(\mathrm{var}(v),t,\mathrm{none}),\\
  \neg\mathrm{occurs}(\mathrm{var}(v),t)
    &\Rightarrow \mathrm{unify\_out}(\mathrm{var}(v),t,
       \mathrm{some}(\mathrm{bind}(v,t))).
\end{align*}
The query asks SNARK to find a result:
\begin{lstlisting}
(prove '(mw-unify-out (mw-var v) term ?result)
       :answer '(values ?result))
\end{lstlisting}
The extracted answer is:
\begin{lstlisting}
(ANSWER-IF (MW-OCCURS (MW-VAR V) TERM)
  (VALUES MW-NONE)
  (VALUES (MW-SOME (MW-BIND V TERM))))
\end{lstlisting}
This is the expected program fragment: fail if the occurs check fails;
otherwise return the singleton binding.

\subsection{Compound Case}

For compound terms, the classical algorithm recursively unifies the first
components, substitutes that result through the remaining components, then
unifies the remaining components and composes substitutions.  The SNARK
obligations encode three cases:
\begin{enumerate}
\item the first recursive unification fails;
\item the first succeeds but the second recursive unification fails;
\item both succeed and their substitutions are composed.
\end{enumerate}
The extracted answer is the following nested conditional:
\begin{lstlisting}
(ANSWER-IF (MW-FAIL (MW-UNIFY M M1))
  (VALUES MW-NONE)
  (ANSWER-IF
    (MW-FAIL
      (MW-UNIFY
        (MW-SUBST (MW-SUBST-OF (MW-UNIFY M M1)) N)
        (MW-SUBST (MW-SUBST-OF (MW-UNIFY M M1)) N1)))
    (VALUES MW-NONE)
    (VALUES
      (MW-SOME
        (MW-COMPOSE
          (MW-SUBST-OF (MW-UNIFY M M1))
          (MW-SUBST-OF
            (MW-UNIFY
              (MW-SUBST (MW-SUBST-OF (MW-UNIFY M M1)) N)
              (MW-SUBST (MW-SUBST-OF (MW-UNIFY M M1)) N1))))))))
\end{lstlisting}
This term mirrors the conventional recursive unification skeleton, while
remaining within SNARK's first-order answer-extraction mechanism.

\section{Working Example II: Square Root and Binary Search}

Manna and Waldinger's binary-search example starts from a specification for
approximating the square root of a nonnegative real number within a positive
tolerance.  Their derivation shows that a binary-search step emerges naturally:
given a rougher interval of width \(2\epsilon\), test whether the exact square
root is in the right half; if so, move the lower endpoint right by
\(\epsilon\), otherwise keep it \cite{MannaWaldinger1985IJCAI,MannaWaldinger1987}.
The local SNARK development abstracts away the full real
arithmetic theory and focuses on this constructive branch.

\subsection{Binary-Search Step}

The predicate \texttt{MW-SQRT-WITHIN} denotes that an answer is within the
desired interval.  The predicate \texttt{MW-LE} abstracts the comparison
\((z+\epsilon)^2 \leq r\).  The two obligations are:
\begin{align*}
  (z+\epsilon)^2 \leq r
    &\Rightarrow \mathrm{sqrt\_within}(r,z+\epsilon,\epsilon),\\
  \neg((z+\epsilon)^2 \leq r)
    &\Rightarrow \mathrm{sqrt\_within}(r,z,\epsilon).
\end{align*}
The extracted SNARK answer is:
\begin{lstlisting}
(ANSWER-IF (MW-LE (MW-SQUARE (MW-PLUS Z EPS)) R)
  (VALUES (MW-PLUS Z EPS))
  (VALUES Z))
\end{lstlisting}
This is the binary-search branch in its smallest form.

\subsection{Recursive Skeleton}

The second square-root example adds a base-case test and a recursive call
\texttt{MW-SQRT} at doubled
tolerance:
\begin{lstlisting}
(ANSWER-IF (MW-LARGE-TOLERANCE R EPS)
  (VALUES ZERO)
  (ANSWER-IF
    (MW-LE
      (MW-SQUARE
        (MW-PLUS (MW-SQRT R (MW-DOUBLE EPS)) EPS))
      R)
    (VALUES (MW-PLUS (MW-SQRT R (MW-DOUBLE EPS)) EPS))
    (VALUES (MW-SQRT R (MW-DOUBLE EPS)))))
\end{lstlisting}
This is faithful to the structural idea in the original derivation: use a
larger tolerance recursively to obtain a rough estimate, then decide which half
of the interval contains the desired value.

\section{Limitations}

These examples are deliberately small.  They are working SNARK demonstrations
of answer extraction for classical synthesis patterns, not full reconstructions
of Manna and Waldinger's deductive-tableau derivations.  In particular:
\begin{itemize}
\item the square-root examples use abstract arithmetic predicates rather than a
complete theory of real numbers;
\item the unification examples isolate key branches rather than proving a
single global total-correctness theorem for all expression pairs;
\item induction and termination arguments are represented by recursive
skeleton symbols rather than synthesized from a full well-founded induction
proof.
\end{itemize}
These choices keep the examples runnable and inspectable inside the existing
SNARK example suite while preserving the core answer-extraction phenomenon.

\section{Running the Examples}

After rebuilding SNARK, start the local executable and enter the SNARK user
package:
\begin{lstlisting}
./snark
(in-package :snark-user)
\end{lstlisting}
Then run:
\begin{lstlisting}
(mw-synthesis-examples)
\end{lstlisting}
The result is a list pairing each example name with its extracted answer term.

\appendix

\section{Prompts and Reproducibility Log}

This appendix records the user-facing prompts and the operational command log
used to produce the development.  It does not include private chain-of-thought.
For reproducibility, the private deliberation is replaced by concise reasoning
summaries and concrete commands.

\subsection{User Prompts}

\begin{enumerate}
\item
\begin{lstlisting}
Compile SNARK for my Mac.
\end{lstlisting}
\item
\begin{lstlisting}
SNARK comes with the capability of answer extraction, which can be used for program synthesis. Manna/Waldinger did this for deductively synthesizing, among other things, the unification algorithm, binary search (via a specification for computing the square root), etc. Look these up and redo these examples in SNARK, under examples/synthesis.
\end{lstlisting}
\item
\begin{lstlisting}
Under docs/synthesis, make a TeX document outlining this new development. Include the classical examples as your two working examples. In it, outline the development, but also include the prompts and the chain-of-thought/agent commands you've applied in an appendix. Use LNCS format for this work.
\end{lstlisting}
\end{enumerate}

\subsection{Research Queries and Sources Consulted}

The literature lookup used the following public-search queries:
\begin{lstlisting}
Manna Waldinger deductive synthesis unification algorithm program synthesis
Manna Waldinger binary search square root deductive synthesis
Manna Waldinger deductive synthesis square root binary search specification
"Deductive synthesis of the unification algorithm" pdf
"The origin of a binary-search paradigm" "sqrt" pdf
"Automatic Deductive Synthesis of the Unification Algorithm" SNARK Waldinger pdf
\end{lstlisting}
The most relevant sources were the SRI pages and PDFs for unification and
binary search, the IJCAI short paper on binary search, and the Isabelle theory
page summarizing a Manna--Waldinger-style unification formalization.

\subsection{Concise Reasoning Summary}

\begin{enumerate}
\item Inspect the repository to identify build scripts and existing example
style.
\item Build SNARK with the installed SBCL, because the checked-in 64-bit SBCL
script points at an obsolete local path.
\item Study existing answer-extraction examples, especially the front/last and
list-reversal examples.
\item Recast the two classical synthesis developments as small SNARK
existence proofs whose answers are program fragments.
\item Wire the new examples into both the manual system loader and ASDF example
system.
\item Rebuild SNARK and verify that the saved native image runs the new
examples directly.
\item Write this LNCS-format document under \texttt{docs/synthesis}.
\end{enumerate}

\subsection{Representative Agent Commands}

\begin{lstlisting}
ls
git status --short
rg --files
find . -maxdepth 2 -iname '*readme*' -o -iname '*install*' -o -iname '*build*'

sed -n '1,220p' INSTALL
sed -n '1,220p' README
sed -n '1,240p' compile
sed -n '1,220p' make-snark-sbcl64
sed -n '1,220p' run-snark

command -v sbcl
command -v ccl
command -v ros
uname -m
sbcl --version
./make-snark-sbcl

ls -lh snark compile.out
file snark
tail -n 80 compile.out
rg -n "ERROR|Error|error|WARNING|Warning|warning|failed|Unhandled|debugger invoked" compile.out

./snark --noinform --disable-debugger \
  --eval '(in-package :snark-user)' \
  --eval '(reverse-example :length 2)' \
  --eval '(sb-ext:quit)'

sed -n '1,260p' examples/front-last-example.lisp
sed -n '1,260p' examples/coder-examples.lisp
rg -n "Answer|:answer|answer|values|declare-function|to-lisp|program" docs README examples src
sed -n '500,590p' src/main.lisp

apply_patch
# Added examples/synthesis/README
# Added examples/synthesis/synthesis-examples.lisp
# Updated snark-examples.asd
# Updated snark-system.lisp

./snark --noinform --disable-debugger \
  --eval '(in-package :snark-user)' \
  --eval '(load "examples/synthesis/synthesis-examples.lisp")' \
  --eval '(format t "~S~%" (mw-synthesis-examples))' \
  --eval '(sb-ext:quit)'

./make-snark-sbcl

./snark --noinform --disable-debugger \
  --eval '(in-package :snark-user)' \
  --eval '(format t "~S~%" (mw-synthesis-examples))' \
  --eval '(sb-ext:quit)'

kpsewhich llncs.cls
command -v latexmk
command -v pdflatex
mkdir -p docs/synthesis
\end{lstlisting}

\section{Extracted Answers}

The rebuilt executable produced the following answers for the two headline
working examples.

\subsection{Unification}

\begin{lstlisting}
(UNIFICATION-VARIABLE-CASE
  (ANSWER-IF (MW-OCCURS (MW-VAR V) TERM)
    (VALUES MW-NONE)
    (VALUES (MW-SOME (MW-BIND V TERM)))))

(UNIFICATION-COMPOUND-CASE
  (ANSWER-IF (MW-FAIL (MW-UNIFY M M1))
    (VALUES MW-NONE)
    (ANSWER-IF
      (MW-FAIL
        (MW-UNIFY
          (MW-SUBST (MW-SUBST-OF (MW-UNIFY M M1)) N)
          (MW-SUBST (MW-SUBST-OF (MW-UNIFY M M1)) N1)))
      (VALUES MW-NONE)
      (VALUES
        (MW-SOME
          (MW-COMPOSE
            (MW-SUBST-OF (MW-UNIFY M M1))
            (MW-SUBST-OF
              (MW-UNIFY
                (MW-SUBST (MW-SUBST-OF (MW-UNIFY M M1)) N)
                (MW-SUBST (MW-SUBST-OF (MW-UNIFY M M1)) N1)))))))))
\end{lstlisting}

\subsection{Binary Search for Square Root}

\begin{lstlisting}
(SQUARE-ROOT-BINARY-SEARCH-STEP
  (ANSWER-IF (MW-LE (MW-SQUARE (MW-PLUS Z EPS)) R)
    (VALUES (MW-PLUS Z EPS))
    (VALUES Z)))

(SQUARE-ROOT-BINARY-SEARCH-PROGRAM
  (ANSWER-IF (MW-LARGE-TOLERANCE R EPS)
    (VALUES ZERO)
    (ANSWER-IF
      (MW-LE
        (MW-SQUARE
          (MW-PLUS (MW-SQRT R (MW-DOUBLE EPS)) EPS))
        R)
      (VALUES (MW-PLUS (MW-SQRT R (MW-DOUBLE EPS)) EPS))
      (VALUES (MW-SQRT R (MW-DOUBLE EPS))))))
\end{lstlisting}

\begin{thebibliography}{8}

\bibitem{MannaWaldinger1981}
Manna, Z., Waldinger, R.:
Deductive synthesis of the unification algorithm.
Science of Computer Programming 1(1--2), 5--48 (1981).
\url{https://www.sri.com/publication/artificial-intelligence-pubs/deductive-synthesis-of-the-unification-algorithm/}

\bibitem{Waldinger2021}
Waldinger, R.:
What, Again? Automatic Deductive Synthesis of the Unification Algorithm.
UNIF 2021 (2021).
\url{https://www.sri.com/publication/artificial-intelligence-pubs/what-again-automatic-deductive-synthesis-of-the-unification-algorithm/}

\bibitem{IsabelleUnification}
Coen, M., Slind, K., Krauss, A.:
Isabelle/HOL theory \texttt{HOL-ex/Unification}.
\url{https://isabelle.in.tum.de/library/HOL/HOL-ex/Unification.html}

\bibitem{MannaWaldinger1985IJCAI}
Manna, Z., Waldinger, R.:
The Origin of the Binary-Search Paradigm.
Proceedings of IJCAI 1985, 223--225 (1985).
\url{https://www.ijcai.org/Proceedings/85-1/Papers/041.pdf}

\bibitem{MannaWaldinger1987}
Manna, Z., Waldinger, R.:
The origin of a binary-search paradigm.
Science of Computer Programming 9(1), 37--83 (1987).
\url{https://www.sri.com/publication/artificial-intelligence-pubs/the-origin-of-the-binary-search-paradigm/}

\bibitem{MannaWaldinger1980}
Manna, Z., Waldinger, R.:
A deductive approach to program synthesis.
ACM Transactions on Programming Languages and Systems 2(1), 90--121 (1980).
\url{https://www.sri.com/publication/artificial-intelligence-pubs/a-deductive-approach-to-program-synthesis/}

\bibitem{SnarkGuide}
Stickel, M., Waldinger, R., Chaudhri, V.:
A Guide to SNARK.
\url{https://nilqed.github.io/SNARK/tutorial/tutorial.html}

\end{thebibliography}

\end{document}
